\chapter{Einleitung} 
\pagenumbering{arabic}
Die direkte Übersetzung neuronaler Aktivität in Steuersignale für externe Geräte hat sich in den letzten Jahrzehnten von einem theoretischen Konzept zu einer klinisch demonstrierten Realität entwickelt. Während Hochberg et al.\ 2012 zeigten, dass Personen mit Tetraplegie einen Roboterarm über ein intrakortikales Implantat steuern konnten \cite{hochberg2012_roboticarm}, gelang Meng et al.\ 2016 ein vergleichbares Ergebnis allein auf Basis nicht-invasiver Elektroenzephalographie: Probanden steuerten durch die mentale Vorstellung von Bewegungsabläufen (\textit{Motor Imagery}) einen Roboterarm und konnten Objekte in einem dreidimensionalen Raum greifen und bewegen \cite{meng2016_eeg_roboticarm}. \newline
Solche EEG-basierten Systeme setzen bislang jedoch hochspezialisierte und kostspielige Laborausstattung voraus. Open-Source-Plattformen wie OpenBCI bieten dabei eine vielversprechende Alternative: Das Ultracortex-IV von OpenBCI ist mit einem Anschaffungspreis von ca.\ \euro{440} um ein Vielfaches günstiger als konventionelle Laborsysteme, die je nach Ausstattung Kosten von bis zu \euro{150000} erreichen können \cite{openbci_ultracortex_shop} \cite{appenman2018_eeg_lab}
